\subsection{``Injectivity radii'' of anti-de Sitter 3-manifolds}\label{inj-radii}
Let $\Gamma$ be a discontinuous group for $\mathrm{AdS}^{3}$. In~this subsection, we give a uniform estimate of the pseudo-distance
between the origin and the second closest point of each $\Gamma$-orbit.

We recall that $\Gamma(\subset \grave{\ }G\times\grave{\ }G)$
acts isometrically on $\mathrm{AdS}^{3}(\cong\grave{\ }G)$ by
$(\gamma_{1},\gamma_{2})x=\gamma_{1} x\gamma_{2}^{-1}$
for~$(\gamma_{1},\gamma_{2})\in\Gamma$ and $x\in\grave{\ }G$. We~set
\begin{align}
\label{varepsilon}
\varepsilon_{\Gamma} := \inf_{(\gamma_1,\gamma_2)\in
\Gamma\setminus\{E\}}
\frac{1}{3}\big|\|\gamma_1\|-\|\gamma_2\|\big|.
\end{align}
By the inequality (see, e.g., \cite[Lemma~5.5]{Kan19})
\begin{gather*}
\|(g_1,g_2)x\|\geq \big | \|g_1\|-\|g_2\| \big | - \|x\|\qquad
\text{for}\quad (g_1,g_2) \in G \quad\text{and}\quad x \in \mathrm{AdS}^{3},
\end{gather*}
we get:
\begin{Lemma}\label{inj-rad}
If $\varepsilon_{\Gamma}>0$, then
$\gamma B(\varepsilon_{\Gamma}) \cap B(\varepsilon_{\Gamma})
= \varnothing$
for all $\gamma\in \Gamma\setminus\{E\}$.
\end{Lemma}

\begin{Proposition}\label{strong}
Let $\Gamma$ be a discrete subgroup of $G$ acting properly discontinuously
on $\mathrm{AdS}^{3}$.
Then there exists $g \in G$
satisfying $\varepsilon_{g^{-1}\Gamma g} > 0$.
\end{Proposition}
\begin{Remark}
One sees in the proof below that
the set of such $g$ is dense in $G$.
\end{Remark}
Proposition~\ref{strong} follows obviously
from the proper discontinuity of the $\Gamma$-action and the following lemma:
\begin{Lemma}\label{weak}
For any countable subset $\Gamma$ of
$G$,
there exists
$g\in G$
such that $\|\gamma_1\|\neq\|\gamma_2\|$
for all $\gamma=(\gamma_1,\gamma_2)\in
g^{-1}\Gamma g \setminus \{E\}$.
\end{Lemma}

\begin{proof}[Proof of Lemma~\ref{weak}]
For $\gamma \in \Gamma$,
the map $f_{\gamma}\colon G \rightarrow G$ defined by $g\mapsto g^{-1}\gamma g$
is real analytic.
For the analytic subset
$F = \{(g_1,g_2) \in G
\mid \|g_1\| = \|g_2\| \}$ of $G$, we claim that
the set $f_{\gamma}^{-1}(F)$ is a proper subset of $G$ if $\gamma\neq E$.
For this, we may assume $\gamma_1 \neq \grave{\ }E$ without loss of generality.
Then there exists
$g_1 \in \grave{\ }G$ satisfying $\|g_1^{-1}\gamma_1g_1\| \neq \|
\gamma_1\|$ as one can find $g_{1}$ depending on the three cases where
$\gamma_{1}$ is hyperbolic, parabolic, or elliptic.
Hence $(g_1,\grave{\ }E)\notin f_{\gamma}^{-1}(F)$ if $\|\gamma_{1}\|=\|\gamma_{2}\|$,
and $E\notin f_{\gamma}^{-1}(F)$ if not. Thus $f_{\gamma}^{-1}(F)$ is a proper subset of $G$.

Therefore the analytic set
$f_{\gamma}^{-1}(F)$ has no interior point, and thus so does the countable union
$\bigcup_{\gamma \in \Gamma \setminus \{E\}}
 f_{\gamma}^{-1}(F)$ by the Baire category theorem
(see, e.g., \cite[Theorem~2.2]{functional-analysis}).
Hence there exists an element
$g$ of $G\setminus\bigcup_{\gamma \in \Gamma \setminus \{E\}}
 f_{\gamma}^{-1}(F)$ and we have
$\|\gamma_1\|\neq\|\gamma_2\|$
for all $\gamma=(\gamma_1,\gamma_2)\in
g^{-1}\Gamma g \setminus \{E\}$.
\end{proof}

\section{Proof of Theorem~\ref{mthm1}}\label{proof-linear-indep}
In this section, we prove Theorem~\ref{mthm1}.

More generally, without finitely generated assumption of $\Gamma$,
we study linear independence of~the generalized Poincar\'e series
of the spherical functions $\psi_{m,k}$ of type $(-m,m+k)$
defined in~Section~\ref{preliminary-spherical}.
By choosing $k=3^{j}$ ($j=0,1,2,\ldots$),
we prove:
\begin{Theorem}
\label{unbound}
If $\Gamma$ is a discontinuous group for $\mathrm{AdS}^{3}$ satisfying
the exponential growth condition~\eqref{C}, then
\begin{gather*}
\lim_{m\to\infty}\mathcal{N}_{\Gamma\backslash\mathrm{AdS}^{3}}(\lambda_{m})=\infty.
\end{gather*}
\end{Theorem}
Theorem~\ref{mthm1} is a direct consequence of
Theorem~\ref{unbound} by Remark~\ref{remPoincare}(2).

\begin{Proposition}
\label{mainthm}
Let $\Gamma$ be a discrete subgroup of $G$ acting properly discontinuously
on $\mathrm{AdS}^{3}$
and satisfying the exponential growth condition \eqref{C}.
If $\varepsilon_{\Gamma}>0$,
then there exists a real number $m_{\Gamma}(k)$
$($given explicitly by $\eqref{Poincare_estimate})$
for $k\in \mathbb{N}$
such that
$\{(\operatorname{Re}(\psi_{m,3^{j}}))^{\Gamma}\}_{j=0}^{k-1}
\subset L^2_{\lambda_{m}}(\Gamma\backslash\mathrm{AdS}^{3})$
are linearly independent
for all integers $m > m_{\Gamma}(k)$.
\end{Proposition}

Postponing the proof of Proposition~\ref{mainthm} until the end of this section,
we prove Theorem~\ref{unbound}.

\begin{proof}[Proof of Theorem~\ref{unbound}]
We have an obvious equality of the multiplicity of $L^2$-eigenvalues,
$\mathcal{N}_{\Gamma\backslash\mathrm{AdS}^{3}} =
\mathcal{N}_{(g^{-1}\Gamma g)\backslash\mathrm{AdS}^{3}}$
for any $g\in G$
through the natural isomorphism
$\Gamma\backslash\mathrm{AdS}^{3} \cong
\big(g^{-1}\Gamma g\big)$ $\backslash\mathrm{AdS}^{3}$ as Lorentzian manifolds.
By replacing $\Gamma$ with $g^{-1}\Gamma g$ if necessary,
we may and do assume $\varepsilon_{\Gamma}>0$ by Proposition~\ref{strong}.
Then Proposition~\ref{mainthm} implies that
$L^2_{\lambda_{m}}\big(\Gamma\backslash \mathrm{AdS}^{3}\big)$ contains at least
$k$ linearly independent elements
if $m>m_{\Gamma}(k)$ for any fixed $k\in\mathbb{N}$,
which means
$\dim_{\mathbb{C}} L^2_{\lambda_{m}}(\Gamma\backslash\mathrm{AdS}^{3})
\geq k.$
Hence Theorem~\ref{unbound} follows.
\end{proof}

Kassel--Kobayashi~\cite{KaKob16} proved the non-vanishing of
the generalized Poincar\'e series $(\psi_{m,0})^{\Gamma}$
for sufficiently large $m\in\mathbb{N}$
by showing that the first term in the generalized Poincar\'e series
is larger at the origin than the sum of the remaining terms.
For this, they utilized the fact that $\psi_{m,0}(\grave{\ }E)=1$.
Our strategy for the proof of Proposition~\ref{mainthm} is along the same line,
however, there are some technical difficulties
since $\psi_{m,k}$ for $k\geq1$ vanishes at the origin. We~then make use of an observation that
$\psi_{m,k}$ decays more slowly at the origin than at infinity,
to be precise, by the following formula, see (\ref{polarFJ}):
\begin{gather*}
|\psi_{m,k}(x)|=\cosh^{-2m}(\|x\|/2) \tanh^{k}(\|x\|/2).
\end{gather*}
Actually, we use an analytic lemma (Lemma~\ref{koukou})
to prove that the first term in the generalized Poincar\'e series $(\psi_{m,k})^{\Gamma}$
is larger at points sufficiently close to the origin
than the sum of the remaining terms if $m\gg0$.
Moreover,
we use a combinatorial lemma (Lemma~\ref{sekoi})
to find points at which
leading terms of
$(\operatorname{Re}(\psi_{m,k}))^{\Gamma}$
do not cancel each other for any linear combination.

For $C,a,\varepsilon > 0$ and $s \in \mathbb{N}$, we set
\begin{align*}%\label{inequality}
m(C,a,\varepsilon,s) := \frac{(\log2)s + 2a\varepsilon + \log\big(1+2^{s}C{\rm e}^{6a\varepsilon}\big)}{\log\cosh\varepsilon}
\end{align*}
and
\begin{align*}
\tilde{m}(C,a,\delta,s) :=
\inf_{0<\varepsilon<\delta}m(C,a,\varepsilon,s).
\end{align*}
Note that $\tilde{m}(C,a,\delta,s)=O\big(\delta^{-2}\big)$ as $\delta\to 0$
and $=O(1)$ as $\delta\to \infty$.
\begin{Lemma}\label{koukou}
For any integer $m > m(C,a,\varepsilon,s)$
and any one-variable polynomial $f$ of degree $\leq s$ with non-negative coefficients,
\begin{gather*}
C\sum_{n=1}^{\infty}{\rm e}^{4a(n+1)\varepsilon}(\cosh 2n\varepsilon)^{-m}f(\tanh
2(n+1)\varepsilon) < (\cosh \varepsilon)^{-m}f(\tanh \varepsilon).
\end{gather*}
\end{Lemma}
\begin{proof}
We may assume that $f(x) = x^{j}$ for $j=0,1,\ldots,s$.
Since
\begin{gather*}
1\leq \frac{\tanh nx}{\tanh x} \leq n,\
(\cosh x)^{n} \leq \cosh nx
\end{gather*}
for $x\in\mathbb{R}$, we have
\begin{align*}
\text{(LHS)/(RHS)} &=
C\sum_{n=1}^{\infty}{\rm e}^{4a(n+1)\varepsilon}\bigg(\frac{\cosh 2n
\varepsilon}{\cosh \varepsilon}\bigg)^{-m}\bigg(\frac{\tanh
2(n+1)\varepsilon}{\tanh \varepsilon}\bigg)^j \nonumber
\\
&\leq C{\rm e}^{6a\varepsilon}\sum_{n=1}^{\infty}
\big({\rm e}^{2a\varepsilon}(\cosh \varepsilon)^{-m}\big)^{2n-1}(2(n+1))^s.
\end{align*}
We set $d:={\rm e}^{2a\varepsilon}(\cosh \varepsilon)^{-m}$.
Then $d<1$ by $m>m(C,a,\varepsilon,s)$. Since $n+1\leq 2^{n}$ for all $n\in\mathbb{N}$,
we have
\begin{gather*}
\text{(LHS)/(RHS)}\leq
2^{s}C{\rm e}^{6a\varepsilon}\sum_{n=1}^{\infty} (2^{s}d)^n
= 2^{s}C{\rm e}^{6a\varepsilon}\frac{2^{s}d}{1-2^{s}d}.
\end{gather*}
Again by $m > m(C,a,\varepsilon,s)$,
we have $2^{s}d<\big(1+2^{s}C{\rm e}^{6a\varepsilon}\big)^{-1}$.
Therefore we obtain
\begin{gather*}
\text{(LHS)/(RHS)}<1.\tag*{\qed}
\end{gather*}
\renewcommand{\qed}{}
\end{proof}
Let $\chi\colon\{\pm 1\} \rightarrow \{0,1\}$ be
the map defined by $\chi(1) = 0$ and $\chi(-1) = 1$.
For $a = (a_j)_{j=0}^{k-1} \in \{\pm1\}^{k}$ and an odd integer $N\geq 3$,
we set
\begin{align*}%\label{theta}
\theta_{a,N} := \pi\sum_{i=0}^{k-1} (\chi(a_i) - \chi(a_{i-1}))N^{-i}.
\end{align*}
Here we use the convention $a_{-1}=1$.
\begin{Lemma}\label{sekoi}
For any $a=(a_{0},\ldots,a_{k-1})\in \{\pm1\}^{k}$ and any odd integer $N$,
we have
\begin{gather*}
a_j\cos\big(N^{j}\theta_{a,N}\big) > 0 \qquad \text{for}\quad j=0,1,\ldots,k-1.
\end{gather*}
\end{Lemma}

\begin{proof}
Since
$N^{k-1}\theta_{a,N} \equiv \pi \chi(a_{k-1}) \ (\mathrm{mod}\ 2\pi)$,
we have $\cos \big(N^{k-1} \theta_{a,N}\big) = a_{k-1}$.
It is easy to check that
$|N^{j} \theta_{a,N} - N^{j} \theta_{(a_0,\cdots,a_j),N}| < \pi/2$ for $j=0,1,\ldots,k-1$,
hence the signature of~$\cos (N^{j} \theta_{a,N})$ is equal to that of
$\cos (N^{j} \theta_{(a_0,\cdots,a_j),N}) = a_j$.
\end{proof}

\begin{Remark}
We have used the geometric progression $(N^{j})_{j=0}^{k-1}$ in Lemma~\ref{sekoi}.
On the other hand, an analogous statement does not hold if
we use arithmetic progressions. For example,
there does not exist $\theta \in \mathbb{R}$
satisfying $a_j\cos m_j\theta > 0$ for all $j=0,1,2,3,4$
if we choose $(a_j)_{j=0}^{4}=(1,1,1,-1,1)$ and
an arithmetic progression $(m_j)_{j=0}^{4}$.
\end{Remark}

For a discontinuous group $\Gamma$ and $k\in \mathbb{N}$,
one can take $m_{\Gamma}(k)$ in Proposition~\ref{Poincare_estimate} by
\begin{gather}
\label{Poincare_estimate}
m_{\Gamma}(k)=\inf_{(A,a)\in \mathrm{C_{exp}}(\Gamma)}
\max\big\{\tilde{m}\big(3^{k-1}A,a,\varepsilon_{\Gamma}/4,3^{k-1}\big)/2,a\big\},
\end{gather}
where $\mathrm{C_{exp}}(\Gamma):=\big\{(A,a)\in \mathbb{R}^{2} \mid
\forall x\in \mathrm{AdS}^{3},\forall R>0, N_{\Gamma}(x,R) < A{\rm e}^{aR}\big\}$.
Here,
we adopt the convention that $\inf_{\varnothing} f =\infty$ for a real-valued function $f$. In~particular,
$m_{\Gamma}(k)=\infty$ when $\mathrm{C_{exp}}(\Gamma)=\varnothing$
or $\varepsilon_{\Gamma}=0$.
\begin{proof}[Proof of Proposition~\ref{mainthm}]
By the exponential growth condition (\ref{C}), $\mathrm{C_{exp}}(\Gamma)\neq\varnothing$
and thus $m_{\Gamma}(k)<\infty$. We~take an integer $m >m_{\Gamma}(k)$.
Then there exist $\varepsilon$ with $0<\varepsilon<\varepsilon_{\Gamma}/4$
and $(A,a)\in\mathrm{C_{exp}}(\Gamma)$
satisfying the inequality
$m > \max\big\{m\big(3^{k-1}A,a,\varepsilon,3^{k-1}\big)/2,a\big\}$.

To see $\mathbb{C}$-linear independence of the real-valued functions
$\big\{(\operatorname{Re}(\psi_{m,3^{j}}))^{\Gamma}\big\}_{j=0}^{k-1}$,
it is enough to prove the non-vanishing of
the real part $\operatorname{Re}\big(\psi_{m,b}^{\Gamma}\big)
=(\operatorname{Re}(\psi_{m,b}))^{\Gamma}$ of
the generalized Poincar\'e series of a linear combination
\begin{gather*}
\psi_{m,b} := \sum_{j = 0}^{k-1} b_j\psi_{m,3^{j}}
\end{gather*}
for any
$b=(b_{0},b_{1},\ldots,b_{k-1}) \in \mathbb{R}^{k}\setminus\{0\}$.
By Lemma~\ref{inj-rad}, for $x\in B(4\varepsilon)$, we have
\begin{gather}
\label{separate}
\psi_{m,b}^{\Gamma}(\Gamma x)
= \psi_{m,b}(x) + \sum_{\substack{\gamma \in \Gamma \\
\|\gamma^{-1} x\| > 4\varepsilon}} \psi_{m,b}\big(\gamma^{-1} x\big).
\end{gather}
By (\ref{polarFJ}), for any $y \in \mathrm{AdS}^{3}$, we get
\begin{gather*}
|\psi_{m,b}(y)| \leq \bigg( \!\cosh \frac{\| y \|}{2} \bigg)^{-2m}
\sum_{j=0}^{k-1} |b_{j}|\bigg(\! \tanh \frac{\| y \|}{2} \bigg)^{3^{j}}.
\end{gather*}
We define
$a=(a_{j})_{j=0}^{k-1}$ by $a_{j}=1$ for $b_{j}\geq 0$ and
$a_{j}=-1$ for $b_{j}<0$,
and set
\begin{gather*}
%\label{pol-b}
f_{b}(u) := \sum_{j=0}^{k-1} b_{j}\cos\big(3^{j}\theta_{a,3}\big)u^{3^{j}}.
\end{gather*}
We note that all the coefficients of $f_{b}$ are non-negative by Lemma~\ref{sekoi}.
Moreover, we get
$\big| \cos \big(3^{j}\theta_{a,3}\big) \big|^{-1}
\leq 3^{k-1}$ for all $j=0,1,\ldots,k-1$
by using the inequality $\sin(\pi x/2)\geq x$ for~$0\leq x\leq1$.
Thus
\begin{gather*}
|\psi_{m,b}(y)| \leq 3^{k-1}\bigg(\!\cosh \frac{\| y \|}{2} \bigg)^{-2m}
f_{b}\bigg(\! \tanh \frac{\| y \|}{2} \bigg)
\end{gather*}
and, for any $x\in B(4\varepsilon)$, we have
\begin{align}
\bigg|\!\!\!\sum_{\substack{\gamma \in \Gamma \\ \|\gamma^{-1} x\| > 4\varepsilon}} \operatorname{Re}\big(\psi_{m,b}\big(\gamma^{-1} x\big)\big)\bigg|
&\leq \sum_{n = 1}^{\infty} \sum_{\substack{\gamma \in \Gamma \\
4\varepsilon n < \|\gamma^{-1} x\| \leq 4\varepsilon (n + 1) }}
|\psi_{m,b}(\gamma^{-1} x)| \nonumber \\
&\leq 3^{k-1}\sum_{n = 1}^{\infty} N_{\Gamma}(x, 4\varepsilon (n + 1))\left(\cosh
2\varepsilon n \right)^{-2m} f_{b}\left(\tanh 2\varepsilon (n+1)\right)
\nonumber \\
&\leq 3^{k-1}A\sum_{n = 1}^{\infty} {\rm e}^{4a\varepsilon (n+1)}\left(\cosh 2\varepsilon n
\right)^{-2m} f_{b}\left(\tanh 2\varepsilon (n+1)\right)
\nonumber \\
&< \left(\cosh \varepsilon\right)^{-2m} f_{b}
\left(\tanh \varepsilon\right).
\label{small}
\end{align}
The third and forth inequalities respectively follow from
the exponential growth condition (\ref{C}) and Lemma~\ref{koukou}.
On the other hand, we set
\begin{gather*}
%\label{non-zero-point}
x_{a,\varepsilon}:=
k\bigg(\frac{\theta_{a,3}}{2}\bigg)a(\varepsilon)k\bigg(\frac{\theta_{a,3}}{2}\bigg)^{-1} \in B(4\varepsilon).
\end{gather*}
Then it follows from (\ref{polarFJ}) that
\begin{align}
\label{large}
\operatorname{Re} \psi_{m,b}(x_{a,\varepsilon})
=\left(\cosh \varepsilon\right)^{-2m} f_{b}
\left(\tanh \varepsilon\right).
\end{align}
By (\ref{separate}), (\ref{small}), and (\ref{large}),
we obtain $(\operatorname{Re}(\psi_{m,b}))^{\Gamma}(\Gamma x_{a,\varepsilon})\neq 0$.
Hence we complete the proof
by the continuity of $\psi_{m,b}^{\Gamma}$ (Fact~\ref{Poincare}).
\end{proof}
